\chapter*{Capaian Pembelajaran dan Peta Kurikulum}
\addcontentsline{toc}{chapter}{Capaian Pembelajaran dan Peta Kurikulum}
\markboth{Capaian Pembelajaran dan Peta Kurikulum}{Capaian Pembelajaran dan Peta Kurikulum}

Bagian ini merangkum capaian pembelajaran, peta Sub-CPMK, dan
keterkaitan setiap bab dengan pertemuan pada RPS. Tabel-tabel di sini
menjadi acuan untuk membaca seluruh buku dan memeriksa ketercapaian
kompetensi.

\section*{Capaian Pembelajaran Lulusan (CPL)}
\addcontentsline{toc}{section}{Capaian Pembelajaran Lulusan (CPL)}

\begin{table}[H]
\centering
\caption{CPL yang didukung mata kuliah}
\label{tab:cpl}
\begin{tabularx}{\textwidth}{@{}lX@{}}
\toprule
\textbf{Kode} & \textbf{Rumusan} \\
\midrule
CPL 02 & Memiliki pengetahuan yang memadai terkait dengan cara kerja
sistem komputer dan mampu merancang dan mengembangkan berbagai
algoritma metode untuk memecahkan masalah. \\
CPL 03 & Kemampuan merancang dan menganalisis algoritma untuk
menyelesaikan permasalahan organisasi secara optimal, serta memilih
dan menerapkannya pada bahasa pemrograman tertentu. \\
\bottomrule
\end{tabularx}
\end{table}

\section*{Capaian Pembelajaran Mata Kuliah (CPMK)}
\addcontentsline{toc}{section}{Capaian Pembelajaran Mata Kuliah (CPMK)}

\begin{table}[H]
\centering
\caption{CPMK dan fokus kompetensi}
\label{tab:cpmk}
\begin{tabularx}{\textwidth}{@{}lX@{}}
\toprule
\textbf{Kode} & \textbf{Rumusan} \\
\midrule
CPMK 01 & Menjelaskan tahapan kompilasi (lexical, syntax, semantic
analysis). \\
CPMK 02 & Memahami konsep Otomata Hingga (Finite Automata). \\
CPMK 03 & Merancang tata bahasa (grammar) bebas konteks untuk bahasa
sederhana. \\
CPMK 04 & Mengimplementasikan parser sederhana untuk validasi kode. \\
\bottomrule
\end{tabularx}
\end{table}

\section*{Peta Sub-CPMK}
\addcontentsline{toc}{section}{Peta Sub-CPMK}

{\small
\begin{longtable}{@{}p{1.2cm}p{7.8cm}p{1.6cm}p{3.5cm}@{}}
\caption{Peta sembilan belas Sub-CPMK ke bab dan asesmen}
\label{tab:subcpmk}\\
\toprule
\textbf{Kode} & \textbf{Rumusan} & \textbf{Bab} & \textbf{Asesmen} \\
\midrule
\endfirsthead
\toprule
\textbf{Kode} & \textbf{Rumusan} & \textbf{Bab} & \textbf{Asesmen} \\
\midrule
\endhead
\bottomrule
\endfoot
1.1 & Menjelaskan perbedaan interpreter dan compiler serta peran
assembler, interpreter, dan compiler sebagai translator, dengan
menyebut contoh masukan dan keluaran masing-masing. (CPL 02) & III &
Kuis Konsep, UAS Tulis \\
1.2 & Menjelaskan fase-fase kompilator (analisis leksikal, sintaksis,
semantik, dan pembangkitan kode antara) beserta peran tabel simbol
dan penanganan kesalahan. (CPL 02) & III & Kuis Konsep, UAS Tulis \\
1.3 & Membuat tabel penelusuran fase kompilasi untuk satu program C
sederhana, yang memuat keluaran antara berupa token, pohon sintaks,
dan kode antara. (CPL 02) & III & UTS Praktik \\
1.4 & Membuat tabel pemetaan hierarki Chomsky (Tipe 3--0) ke jenis
automata dan menguji klasifikasi lima contoh grammar ke kelas yang
tepat. (CPL 02) & IV & Kuis Konsep, UTS Praktik \\
1.5 & Membuat tabel aturan semantik berbasis \textit{syntax-directed
definition} untuk ekspresi dan pernyataan subset bahasa C, lalu
menguji keluarannya pada pohon sintaks. (CPL 02) & XIV & Praktikum 4 \\
1.6 & Mengimplementasikan \textit{symbol table} bersarang dan pemeriksa
tipe dalam bahasa C, dengan berkas kode dan keluaran diagnosis sebagai
bukti. (CPL 02) & XV & Praktikum 4 \\
1.7 & Menguji dan men-\textit{debug} penanganan kesalahan semantik
sehingga program menghasilkan pesan kesalahan yang menyebut baris,
nama, dan jenis pelanggaran. (CPL 02) & XVI & Praktikum 4, Proyek
Akhir \\
2.1 & Membuat tabel ekspresi reguler untuk setiap kategori token subset
bahasa C (kata kunci, \textit{identifier}, bilangan, operator, dan
pemisah). (CPL 02) & V & Praktikum 1 \\
2.2 & Mengimplementasikan simulator NFA dan DFA dalam bahasa C untuk
mengenali token, dengan berkas kode dan keluaran pengujian sebagai
bukti. (CPL 02) & VI & Praktikum 1, UTS Praktik \\
2.3 & Membuat tabel konversi NFA ke DFA dan tabel minimalisasi DFA,
serta menguji bahwa keduanya menerima bahasa token yang sama.
(CPL 02) & VI & Praktikum 1, UTS Praktik \\
2.4 & Mengimplementasikan \textit{lexer} subset bahasa C dalam dua
versi, yaitu \textit{hand-written} dan Flex, lalu menguji keluarannya
pada berkas masukan uji. (CPL 02) & VII & Praktikum 1, Proyek Akhir \\
3.1 & Membuat berkas grammar BNF dan EBNF untuk subset bahasa C, serta
menguji derivasi dan pohon sintaks yang dihasilkan. (CPL 03) & VIII &
Praktikum 2, UTS Praktik \\
3.2 & Menguji grammar subset bahasa C dan men-\textit{debug}
ambiguitas, \textit{left recursion}, serta \textit{left factoring}
sampai berkas grammar bebas konflik. (CPL 03) & IX & Praktikum 2 \\
3.3 & Membuat tabel himpunan FIRST dan FOLLOW untuk grammar subset
bahasa C dan menguji keputusan parser LL(1) berdasarkan tabel
tersebut. (CPL 03) & X & Praktikum 2, UTS Praktik \\
3.4 & Membuat berkas grammar bebas konteks untuk subset bahasa C yang
mencakup deklarasi, ekspresi, dan pernyataan, sebagai dasar parser dan
proyek akhir. (CPL 03) & VIII & Praktikum 2, Proyek Akhir \\
4.1 & Mengimplementasikan \textit{recursive descent parser} dan
\textit{predictive parser} LL(1) dalam bahasa C, dengan berkas kode
dan keluaran validasi sebagai bukti. (CPL 03) & X & Praktikum 2 \\
4.2 & Membuat tabel aksi dan \textit{goto} untuk parser SLR/LALR, lalu
menguji dan men-\textit{debug} konflik \textit{shift-reduce} atau
\textit{reduce-reduce} yang muncul. (CPL 03) & XI--XII & Praktikum 3,
UTS Praktik \\
4.3 & Mengimplementasikan parser subset bahasa C dengan Flex dan Bison
beserta aksi semantik, dengan berkas kode dan keluaran program sebagai
bukti. (CPL 03) & XIII & Praktikum 3, Proyek Akhir \\
4.4 & Membuat AST dan validator kode untuk subset bahasa C, lalu
menguji keluaran pesan kesalahan sintaksis yang informatif.
(CPL 03) & XIII, XVI & Praktikum 3, Proyek Akhir \\
\end{longtable}
}

\section*{Peta Pertemuan dengan Bab}
\addcontentsline{toc}{section}{Peta Pertemuan dengan Bab}

{\small
\begin{longtable}{@{}p{0.9cm}p{5.8cm}p{1.6cm}p{2.6cm}p{2.7cm}@{}}
\caption{Peta enam belas pertemuan RPS ke bab}
\label{tab:peta-pertemuan}\\
\toprule
\textbf{Pert.} & \textbf{Topik} & \textbf{Bab} & \textbf{Sub-CPMK} &
\textbf{Asesmen} \\
\midrule
\endfirsthead
\toprule
\textbf{Pert.} & \textbf{Topik} & \textbf{Bab} & \textbf{Sub-CPMK} &
\textbf{Asesmen} \\
\midrule
\endhead
\bottomrule
\endfoot
1 & Pengenalan kompilasi dan arsitektur kompilator & III & 1.1--1.3 &
Kuis Konsep \\
2 & Teori bahasa formal dan hierarki Chomsky & IV & 1.4 &
Kuis Konsep \\
3 & Analisis leksikal: token, lexeme, pattern, regex & V & 2.1 &
Praktikum 1 \\
4 & Finite automata: NFA, DFA, konversi, minimalisasi & VI & 2.2--2.3 &
Praktikum 1 \\
5 & Implementasi lexical analyzer & VII & 2.4 & Praktikum 1 \\
6 & Context-free grammar: BNF, EBNF, derivasi, parse tree & VIII &
3.1, 3.4 & Praktikum 2 \\
7 & Ambiguitas, left recursion, left factoring & IX & 3.2 &
Praktikum 2 \\
8 & Ujian Tengah Semester (UTS) & --- & CPMK 01--03 & UTS Praktik \\
9 & Top-down parsing: recursive descent, LL(1), FIRST/FOLLOW & X &
3.3, 4.1 & Praktikum 2 \\
10 & Bottom-up parsing: shift-reduce, operator precedence & XI & 4.2 &
Praktikum 3 \\
11 & LR parsing: canonical LR, SLR, LALR, konflik & XII & 4.2 &
Praktikum 3 \\
12 & Parser generator dan semantic actions & XIII & 4.3--4.4 &
Praktikum 3 \\
13 & Analisis semantik: tabel aturan semantik (SDD) & XIV & 1.5 &
Praktikum 4 \\
14 & Symbol table, scope management, type checking & XV & 1.6 &
Praktikum 4 \\
15 & AST, validasi kode, penanganan error & XVI & 4.4, 1.7 &
Praktikum 4, Proyek Akhir \\
16 & Ujian Akhir Semester (UAS) & --- & CPMK 01--04 &
UAS Tulis, Proyek Akhir \\
\end{longtable}
}

\section*{Bobot Penilaian}
\addcontentsline{toc}{section}{Bobot Penilaian}

\begin{table}[H]
\centering
\caption{Komponen penilaian dan bobot}
\label{tab:bobot-penilaian}
\begin{tabularx}{\textwidth}{@{}lXr@{}}
\toprule
\textbf{Komponen} & \textbf{Cakupan} & \textbf{Bobot} \\
\midrule
Praktikum 1 & Lexer (Sub-CPMK 2.1--2.4) & 12\% \\
Praktikum 2 & Grammar dan parser (Sub-CPMK 3.1--3.4, 4.1) & 12\% \\
Praktikum 3 & Tabel SLR/LALR, Bison, dan AST (Sub-CPMK 4.2--4.4) & 12\% \\
Praktikum 4 & Symbol table dan analisis semantik (Sub-CPMK 1.5--1.7) & 12\% \\
Proyek Akhir & Front-end mini-compiler subset C & 30\% \\
UTS Praktik & Ujian praktik (Sub-CPMK 1.3, 2.2, 2.3, 3.1, 3.3, 4.2) & 14\% \\
Kuis Konsep & Konsep kompilasi (Sub-CPMK 1.1, 1.2, 1.4) & 3\% \\
UAS Tulis & Konsep dan analisis (Sub-CPMK 1.1, 1.2, 1.4) & 5\% \\
\midrule
\textbf{Total} & & \textbf{100\%} \\
\bottomrule
\end{tabularx}
\end{table}

Komposisi bobot adalah 92\% praktik (Praktikum 1--4, Proyek Akhir, dan
UTS Praktik) dan 8\% teori (Kuis Konsep dan UAS Tulis).
